% ==============================================================================
% dxh_flight_principle_3d.tex — DXH workshop: quadcopter flight principle (3D)
% DXH講座: クアッドコプタ飛行原理の3D図
%
% Isometric-style view of an X-quad: four rotor discs with thrust vectors,
% spin-direction arcs (CW/CCW), and the three body axes with roll/pitch/yaw
% moment arrows. Projection: forward = (-0.78, 0.59), right = (0.98, 0.20),
% up = (0, 1); rotor arm length a = 1.55.
% X型クアッドの斜視図: 4つのロータ円盤+推力ベクトル、回転方向の円弧
% (CW/CCW)、機体3軸とロール/ピッチ/ヨーのモーメント矢印。
% 投影: 前方=(-0.78,0.59), 右=(0.98,0.20), 上=(0,1)、アーム長 a=1.55。
% ==============================================================================
\documentclass[border=10pt]{standalone}
\usepackage{tikz}
\usetikzlibrary{arrows.meta}
\usepackage{luatexja}
\usepackage[match]{luatexja-fontspec}
% Cross-platform font: Hiragino (macOS) → Yu Gothic (Windows) → Noto (Linux)
\usepackage{ifthen}
\newboolean{jfontset}\setboolean{jfontset}{false}
\IfFontExistsTF{Hiragino Kaku Gothic ProN}{%
  \setmainjfont{Hiragino Kaku Gothic ProN}%
  \setsansjfont{Hiragino Kaku Gothic ProN}%
  \setboolean{jfontset}{true}}{}
\ifthenelse{\boolean{jfontset}}{}{%
  \IfFontExistsTF{Yu Gothic}{%
    \setmainjfont{Yu Gothic}%
    \setsansjfont{Yu Gothic}%
    \setboolean{jfontset}{true}}{}}
\ifthenelse{\boolean{jfontset}}{}{%
  \IfFontExistsTF{Noto Sans CJK JP}{%
    \setmainjfont{Noto Sans CJK JP}%
    \setsansjfont{Noto Sans CJK JP}}{}}

\begin{document}

% Colors defined after \begin{document} (standalone compatibility, rule T9)
% 色定義は \begin{document} の後（standalone 互換, ルール T9）
\definecolor{cthrust}{RGB}{180, 60, 40}   % thrust vectors / 推力
\definecolor{ccw}{RGB}{0, 100, 180}       % CW rotor spin / CW回転
\definecolor{cccw}{RGB}{200, 90, 40}      % CCW rotor spin / CCW回転
\definecolor{cmom}{RGB}{40, 140, 60}      % moment arrows / モーメント

\begin{tikzpicture}[
    lbl/.style={font=\small\bfseries, align=center},
    axis/.style={dashed, black!45, line width=0.8pt},
    thrust/.style={-{Stealth[length=3mm]}, cthrust, line width=1.6pt},
    mom/.style={-{Stealth[length=2.6mm]}, cmom, line width=1.4pt},
    spin/.style={-{Stealth[length=2mm]}, line width=1.1pt},
  ]

  % Rotor centers (isometric projection) / ロータ中心（斜視投影）
  % FL=(-2.17,0.31)  FR=(0.62,0.87)  RL=(-0.62,-0.87)  RR=(2.17,-0.31)

  % === Body axes (behind everything) / 機体3軸（最背面） ===
  \draw[axis] (0,0) -- (-3.05,2.3);            % forward axis / 前後軸
  \draw[axis] (0,0) -- (3.15,0.64);            % lateral axis / 左右軸
  \draw[axis] (0,0) -- (0,-2.05);              % vertical axis / 上下軸

  % === Arms / アーム ===
  \draw[line width=2.4pt, black!65] (0,0) -- (-2.17,0.31);
  \draw[line width=2.4pt, black!65] (0,0) -- (0.62,0.87);
  \draw[line width=2.4pt, black!65] (0,0) -- (-0.62,-0.87);
  \draw[line width=2.4pt, black!65] (0,0) -- (2.17,-0.31);

  % === Rotor discs / ロータ円盤 ===
  \draw[fill=black!12, black!60] (-0.62,-0.87) ellipse [x radius=0.62, y radius=0.24];
  \draw[fill=black!12, black!60] (2.17,-0.31) ellipse [x radius=0.62, y radius=0.24];
  \draw[fill=black!12, black!60] (-2.17,0.31) ellipse [x radius=0.62, y radius=0.24];
  \draw[fill=black!12, black!60] (0.62,0.87) ellipse [x radius=0.62, y radius=0.24];

  % === Center body + nose / 機体中心と機首 ===
  \draw[fill=black!30, black!60] (0,0) ellipse [x radius=0.34, y radius=0.16];
  \node[lbl, text=black!70] at (-1.35,1.42) {前};

  % === Spin direction arcs / 回転方向の円弧 ===
  % CW rotors: FL (M4), RR (M2) — blue / CWロータ: FL・RR（青）
  \draw[spin, ccw] ([shift={(-2.17,0.31)}] 210:0.8 and 0.34)
    arc[x radius=0.8, y radius=0.34, start angle=210, end angle=-30];
  \draw[spin, ccw] ([shift={(2.17,-0.31)}] 210:0.8 and 0.34)
    arc[x radius=0.8, y radius=0.34, start angle=210, end angle=-30];
  % CCW rotors: FR (M1), RL (M3) — orange / CCWロータ: FR・RL（橙）
  \draw[spin, cccw] ([shift={(0.62,0.87)}] -30:0.8 and 0.34)
    arc[x radius=0.8, y radius=0.34, start angle=-30, end angle=210];
  \draw[spin, cccw] ([shift={(-0.62,-0.87)}] -30:0.8 and 0.34)
    arc[x radius=0.8, y radius=0.34, start angle=-30, end angle=210];
  \node[font=\scriptsize\bfseries, text=ccw]  at (3.05,-0.75) {CW};
  \node[font=\scriptsize\bfseries, text=cccw] at (-1.9,-1.35) {CCW};

  % === Thrust vectors / 推力ベクトル ===
  \draw[thrust] (-2.17,0.58) -- ++(0,1.15);
  \draw[thrust] (0.62,1.14) -- ++(0,1.15);
  \draw[thrust] (-0.62,-0.60) -- ++(0,1.15);
  \draw[thrust] (2.17,-0.04) -- ++(0,1.15);
  \node[lbl, text=cthrust] at (1.35,2.35) {推力};

  % === Moment arrows around each axis / 各軸まわりのモーメント矢印 ===
  % Roll (around forward axis) / ロール（前後軸まわり）
  \begin{scope}[rotate around={143:(-2.85,2.15)}]
    \draw[mom] ([shift={(-2.85,2.15)}] 185:0.38 and 0.5)
      arc[x radius=0.38, y radius=0.5, start angle=185, end angle=-40];
  \end{scope}
  \node[lbl, text=cmom] at (-3.55,1.45) {ロール};

  % Pitch (around lateral axis) / ピッチ（左右軸まわり）
  \begin{scope}[rotate around={11:(2.95,0.60)}]
    \draw[mom] ([shift={(2.95,0.60)}] 185:0.38 and 0.5)
      arc[x radius=0.38, y radius=0.5, start angle=185, end angle=-40];
  \end{scope}
  \node[lbl, text=cmom] at (3.6,-0.05) {ピッチ};

  % Yaw (around vertical axis) / ヨー（上下軸まわり）
  \draw[mom] ([shift={(0,-1.75)}] 200:0.6 and 0.24)
    arc[x radius=0.6, y radius=0.24, start angle=200, end angle=-30];
  \node[lbl, text=cmom] at (1.05,-1.85) {ヨー};

\end{tikzpicture}

\end{document}
